\section{Analysis}
\label{sec:theory}

This section analyzes the shared-factor Tucker polynomial and its estimator. It
first connects the general $p$-fold model to NTD-PL, then studies the resulting
rank and parameter cost. The remaining results cover response approximation,
descent, and computational scaling. The optimization results also apply to
other differentiable finite response bases whose response is linear in the
coefficients.

\subsection{Tucker Polynomial Structure and Rank Inflation}

We begin with the effect of sharing the mode factors across polynomial orders.

\begin{proposition}[Shared-factor $p$-fold Tucker term]
\label{prop:shared_pfold}
Let $\cS$ be the Tucker tensor in \eqref{eq:latent_tucker}. If the mode maps of
the degree-$p$ term satisfy \eqref{eq:shared_pfold_factors}, then
\begin{equation}
  \mathcal{T}_p(\cX)=\cS^{\od p}.
\end{equation}
Consequently, the shared-factor Tucker polynomial is exactly NTD-PL.
\end{proposition}

This result explains why the two views of the model agree. The general
$p$-fold term is a polynomial map of the core. Sharing its mode factors turns
that term into an entrywise power of one Tucker tensor.

\begin{proposition}[Response-induced Tucker rank inflation]
\label{prop:rank_inflation}
Let $\cS$ have Tucker rank at most $(R_1,\ldots,R_N)$. For every $p \ge 1$,
$\cS^{\od p}$ admits a Tucker representation with mode ranks at most
\begin{equation}
  \left(
  \binom{R_1+p-1}{p},\ldots,\binom{R_N+p-1}{p}
  \right).
\end{equation}
\end{proposition}

The quantity $\binom{R_n+p-1}{p}$ is the number of degree-$p$ monomials in a
row of the mode-$n$ factor. Thus each shared $p$-fold term has an ordinary
Tucker representation with polynomially expanded factors, even though those
factors are determined by the original ones.

\begin{corollary}[Parameter cost of sharing]
\label{cor:tied_lift_parameters}
The direct degree-$p$ term in \eqref{eq:pfold_tucker} contains
$\sum_{n=1}^{N}I_nR_n^p$ mode-map entries. Let
$D_n(p)=\binom{R_n+p-1}{p}$. An ordinary Tucker model that explicitly
represents the shared degree-$p$ term may require
\begin{equation}
  d_{\mathrm{lift}}(p)=\prod_{n=1}^{N}D_n(p)+\sum_{n=1}^{N}I_nD_n(p)
\end{equation}
parameters for that term alone. The shared-factor model represents all degrees
$0,\ldots,P$ with the original Tucker parameters plus $P+1$ response
coefficients.
\end{corollary}

Sharing therefore reduces both constructions: it avoids separate $p$-fold mode
maps and avoids storing the larger Tucker representation induced by the
polynomial terms. This economy applies only to interactions generated from the
same core and factors; the general Tucker polynomial remains more flexible.

\begin{theorem}[Separation from matched-rank Tucker]
\label{thm:tucker_separation}
Let \(N\ge 3\), \(R\ge 2\), and \(p\ge 2\). Define
\(D=\binom{R+p-1}{p}\) and suppose \(I_n\ge D\) for every mode. Then there
exists an order-\(N\) tensor representable by a degree-\(p\)
NTD-PL model with Tucker rank \((R,\ldots,R)\), but whose
exact standard Tucker rank is \((D,\ldots,D)\).
\end{theorem}

The theorem shows that sharing does not reduce the model to matched-rank
Tucker. A shared $p$-fold term can attain the expanded ranks in
Proposition~\ref{prop:rank_inflation}, while the trainable Tucker tensor remains
at rank \((R,\ldots,R)\). This is an existence result, not a claim that NTD-PL
spans all tensors of the larger rank.

\subsection{Response Generation and Approximation}

The previous results explain how a polynomial response can inflate the Tucker
rank observed after measurement. The next two results ask what kinds of
entrywise responses the power basis can represent on top of the same
low-rank Tucker signal.

\begin{proposition}[Exact generation for polynomial responses]
\label{prop:poly_generation}
If $\cY=g(\cS^\star)$, where $\cS^\star$ has Tucker rank at most
$(R_1,\ldots,R_N)$ and $g$ is a polynomial of degree at most $P$, then a
degree-$P$ NTD-PL model with the same Tucker rank can
represent $\cY$ exactly.
\end{proposition}

\begin{theorem}[Continuous response approximation]
\label{thm:response_approximation}
Let $\cS^\star$ have Tucker rank at most $(R_1,\ldots,R_N)$ and entries in a
compact interval $[a,b]$. If the measurement response $f$ is continuous on
$[a,b]$, then for every $\epsilon>0$ there exists a polynomial degree $P$ and
coefficients $\bbeta$ such that
\begin{equation}
  \norm{f(\cS^\star)-f_{\bbeta}(\cS^\star)}_F
  \le
  \epsilon \sqrt{M},
  \qquad
  M=\prod_{n=1}^{N}I_n .
\end{equation}
\end{theorem}

The exact polynomial result is a useful special case. The approximation result
is broader: the power response is not restricted to exactly polynomial sensors.
On a bounded signal range, it can approximate any smooth or continuous
entrywise response to the desired uniform accuracy, while keeping the Tucker
rank of the underlying signal fixed. The degree required for a given accuracy
depends on the regularity of the response and the width of the signal range.
This approximation result is the reason the power basis is used as the
default response basis in the experiments.

\subsection{Descent and Stationarity}

We analyze the fixed-basis objective $\mathcal{F}_J$ in
\eqref{eq:estimation_objective}. For nested bases, the continuation schedule
can be viewed as a sequence of such fixed-basis problems, each initialized from
the previous basis size.

\begin{assumption}
\label{ass:optimization}
For a fixed active basis size $J+1$, the iterates remain in a bounded normalized
parameter set. The regularization weights are nonnegative and
$\lambda_\beta>0$. The Tucker step either uses a line search or an equivalent
descent rule such that, for constants $c>0$ and $b>0$,
\begin{align}
  \mathcal{F}_J(\theta^k)-\mathcal{F}_J(\theta^{k+1})
  &\ge c\norm{\theta^{k+1}-\theta^k}_2^2,
  \label{eq:sufficient_decrease}\\
  \mathrm{dist}\!\left(0,\partial \mathcal{F}_J(\theta^{k+1})\right)
  &\le b\norm{\theta^{k+1}-\theta^k}_2 .
  \label{eq:relative_error}
\end{align}
\end{assumption}

The boundedness condition is enforced in practice by Tucker gauge
normalization and regularization. The two descent conditions are standard for
block coordinate and line-search gradient schemes on smooth objectives.

\begin{proposition}[Monotone objective decrease]
\label{prop:monotone_descent}
Under Assumption~\ref{ass:optimization}, every fixed-basis outer iteration of
the alternating estimator satisfies
\begin{equation}
  \mathcal{F}_J(\theta^{k+1})\le \mathcal{F}_J(\theta^k).
\end{equation}
Moreover, if \eqref{eq:sufficient_decrease} holds for the entire outer update,
then
\begin{equation}
  \sum_{k=0}^{\infty}
  \norm{\theta^{k+1}-\theta^k}_2^2
  <\infty .
\end{equation}
\end{proposition}

\begin{theorem}[Convergence to a stationary point]
\label{thm:stationary_convergence}
Suppose Assumption~\ref{ass:optimization} holds for a fixed basis size
$J+1$. Then
the sequence generated by the alternating estimator has finite length and
converges to a critical point $\theta^\star$ of $\mathcal{F}_J$:
\begin{equation}
  0\in \partial \mathcal{F}_J(\theta^\star).
\end{equation}
\end{theorem}

The theorem does not claim global optimality; the problem remains nonconvex
because of the Tucker factorization. It applies to the descent-controlled
implementation specified in Assumption~\ref{ass:optimization}. The Adam solver
used in the experiments is a practical implementation and is not claimed to
satisfy these conditions at every iteration. Under the stated conditions, the
result places the alternating scheme on the same theoretical footing as standard
nonconvex block-descent methods
\cite{Attouch2013Convergence,Bolte2014PALM}.

\subsection{Computational Scaling}

Let $M=\prod_{n=1}^{N}I_n$, $R_\Pi=\prod_{n=1}^{N}R_n$, and
$d_{\mathrm{Tucker}}=R_\Pi+\sum_{n=1}^{N}I_nR_n$. Storing the fitted model
requires
\begin{equation}
  d_{\mathrm{NTD}}=d_{\mathrm{Tucker}}+(J+1)
\end{equation}
parameters. Thus the learned response adds only $J+1$ parameters to the Tucker
model. For one outer iteration, Response Refresh forms
$\bPhi_{\Omega,J}^{T}\bPhi_{\Omega,J}$ and solves a small system, costing
\begin{equation}
  \mathcal{O}\!\left(mJ^2+J^3\right).
\end{equation}
The dominant cost remains the Tucker reconstruction and reverse contractions;
evaluating the response adds $\mathcal{O}(MJ)$ work for dense tensors or
$\mathcal{O}(mJ)$ on observed entries. Detailed scaling comparisons, the
projection interpretation of $\Dlink$, and a masked-entry generalization bound
are provided in the supplementary material.
